% Source presentations for register row N-c-006.

\subsection{The two-variable decision procedure of Gr\"adel--Kolaitis--Vardi}
\label{sec:p09-presentation-n-c-006}

The language \(\mathrm{FO}^2\) in this presentation is relational first-order
logic whose formulas use at most two distinct individual variables, repeatedly
rebound as needed.  Equality is admitted unless the equality-free variant is
explicitly named.  Relation symbols may initially have higher arity: because
only two variables occur, each atomic occurrence has one of finitely many
variable patterns, and the source replaces those patterns by fresh unary or
binary relation symbols before the model construction
\cite[pp.~53--58]{GradelKolaitisVardi1997}.

Proposition~3.1 gives a satisfiability-preserving Scott reduction.  From a
sentence \(\varphi\) of size \(s\), it constructs a sentence \(\varphi^*\) of
size \(O(s\log s)\) using \(O(s)\) relation symbols and having the form
\[
 \forall x\forall y\,\alpha(x,y)
 \ \land\
 \bigwedge_{i=1}^{m}\forall x\exists y\,\beta_i(x,y),
\]
where the matrices are quantifier-free and may be arranged so that each
\(\beta_i(x,y)\) entails \(x\neq y\).  The displayed conjunction is the
source's \emph{Scott class}.  It is not one prenex sentence, and calling it a
prefix class would erase the alternation among its conjuncts
\cite[prop.~3.1 and pp.~58--60]{GradelKolaitisVardi1997}.

A \(k\)-type is a maximal consistent set of atomic and negated atomic formulas
in a fixed \(k\)-tuple of variables.  An element is a \emph{king} when it is
the unique realizer of its one-type.  These definitions are needed because the
finite-model construction must preserve unique realizers exactly while it may
copy nonunique types.  The set of kings together with selected witnesses for
their existential obligations is the \emph{royal court}; this phrase denotes
that explicitly constructed finite set and nothing beyond the proof
\cite[defs.~4.1--4.2]{GradelKolaitisVardi1997}.

If a Scott-class sentence \(\theta\) has size \(s\) and uses \(r\) relation
symbols, Theorem~4.3 constructs a finite model of size at most
\[
 3s\,2^r.
\]
Starting from any model, retain every king, close the royal court under one
chosen witness for each \(\forall x\exists y\) conjunct, and create three
cyclic banks of representatives for every remaining one-type and witness
index.  Realized two-types from the original model are copied so that the
universal conjunct remains true.  The banks assign witnesses cyclically,
preventing the reciprocal-witness conflict that can occur with only two banks;
unused pairs receive an arbitrary compatible two-type already realized in the
original model.  The source counts the resulting universe by
\[
 (m+1)\lvert K\rvert
 +3m\bigl(\lvert P\rvert-\lvert K\rvert\bigr)
 \leq 3s\,2^r,
\]
where \(K\) is the set of kings and \(P\) the set of realized one-types
\cite[pp.~60--64]{GradelKolaitisVardi1997}.

Combining this construction with the Scott reduction yields the exponential
model property: for some constant \(c\), every satisfiable \(\mathrm{FO}^2\)
sentence of size \(s\) has a model of size at most \(2^{cs}\).  Theorem~5.3
places the Scott class in nondeterministic time
\(2^{O(s/\log s)}\) and arbitrary \(\mathrm{FO}^2\) in
\(2^{O(s)}\).  A matching lower bound already holds without equality, so
satisfiability for \(\mathrm{FO}^2\), with or without equality, is
NEXPTIME-complete \cite[thms.~5.1 and 5.3 and cor.~5.4]{GradelKolaitisVardi1997}.
The finer Scott-class bound is not identical to the class-level completeness
statement because the reduction itself expands the formula by
\(O(s\log s)\).

Constants can be treated as kings without changing the exponential bound.
Positive-arity function symbols are different: even unary functions lead to
undecidable two-variable cases.  The three-variable fragment is already
undecidable without equality.  Counting quantifiers preserve decidability but
not the ordinary finite-model property, and a nearby unrestricted
\(\forall\forall\land\forall\exists\) conjunction can express an infinite
linear order without a maximum \cite[pp.~53--54 and 66--67]{GradelKolaitisVardi1997}.
These are boundary results for extensions of the language, not exceptions to
the relational \(\mathrm{FO}^2\) theorem just reconstructed.

\paragraph{Mortimer 1975 source boundary.}
Mortimer's six-page article is the registered original finite-model result, but
its theorem text was not available for source-specific verification.  The 1997
authors report that Mortimer proved the finite-model property with a doubly
exponential bound; their paper improves that bound by one exponential.  The
Scott conjunction, types, kings, three-bank construction, and improved bound
above are therefore attributed only to the 1997 paper and not retrojected into
Mortimer's presentation \cite{Mortimer1975,GradelKolaitisVardi1997}.
